\section{The Operator Turns}

The concrete operator has now acquired every responsibility the chapter
needed to assign. It reads neighboring differences. It is anchored against
null drift. It carries a positive energy square. Its square root supplies a
scale field. Its invariant binds the labels by which a universe can be read.
Its records are physical or metaphysical only relative to sensor classes
that can or cannot resolve them. What remains is the turn itself: the
derivative pipeline must be instantiated over this operator.

The word instantiated is important. The chapter is not adding a second
grammar of derivatives beside the finite gauge. It is showing that the
derivative grammar from the previous chapter has found a concrete reader.
Residual, first variation, linearization, and second variation are not four
detached notions. They are successive readings of the same finite activity
once the concrete operator supplies the chain, boundary, and scale.

A residual is the unread imbalance left by a proposed reading. On the
anchored chain, a residual appears when the neighboring differences fail to
cancel under the permitted test. If every supported test reads no imbalance,
the residual has vanished for that sensor class. If one supported test
still reads a debt, the proposed reading is not stationary. The operator
turns because it gives the residual a place to be measured rather than a
name to be admired.

The first variation is the residual read linearly. It asks how the activity
changes under a permitted small turn of the chain. Since the chain is
finite, the small turn does not need to be imagined as an infinite
infinitesimal substance. It is a disciplined comparison of nearby permitted
readings, with the leading imbalance extracted by the grammar. The first
variation is therefore the residual in its testable direction.

Linearization is the same turn seen as a rule. Once the leading response has
been read, the grammar can ask whether any other leading rule would give the
same residual readings for all permitted tests. Once the concrete operator
carries those tests, the answer is no. The leading rule is unique because
the supported tests determine it. A different linearization would either
fail to read an available residual or invent one the operator does not
carry.

Stationarity is the moment when that leading rule reads zero. The first
variation vanishes for every permitted test. The derivative has not
disappeared into vagueness; it has resolved to the zero leading imbalance.
At that point the first surviving activity is quadratic. This is the same
second variation that Chapter 02 identified as the unique positive
invariant, but now the invariant is grounded in the concrete operator.

The finite Taylor picture helps if it is kept grammatical. A proposed
reading is varied. The constant term records the old reading. The linear
term is the first residual. The quadratic term is the activity left after
the first residual has vanished. In this finite setting the local story
terminates at order two because the operator has been built to read the
activity square. Higher description may be useful later, but it is not the
first positive invariant. The quadratic remainder is the one the grammar is
forced to carry.

This also explains the difference between telescoping cancellation and
discriminating obstruction. A telescoping chain cancels its interior debts:
each local contribution is answered by the next, and the boundary leaves no
unspent residue. The derivative reading is zero. A discriminating chain
leaves one supported debt that cannot be cancelled by the same telescope.
The sensor class detects nonstationarity. The distinction is not a matter of
style. It is the operator reading whether the residual has a witness.

The concrete operator therefore turns the whole pipeline. Through it, the
burden read by one clock is placed on an anchored chain; its energy is
measured; its energy becomes scale; its labels are bound; its readability is
classified by sensor class; and its derivative content is read. Each step is
forced by the prior grammar. Nothing in the pipeline is allowed to float as
a merely named abstraction.

The chapter's title has now paid its debt. The operator is concrete because
it instantiates the finite invariant in a reader that can detect zero
activity and measure nonzero activity. Metaphysical relativity is present
because the same invariant may be physical for one sensor class and
metaphysical for another. The two phrases belong together. Without the
operator, relativity would be only a slogan about perspectives. Without
metaphysical relativity, the operator would be mistaken for an absolute
view from nowhere.

What the grammar has instead is stricter and more useful. It has one
positive invariant, one concrete finite reader, and many possible sensor
classes. It can say what resolves, what obstructs, what is quotiented, and
what remains unread. The next chapter can now let that grounded invariant
split into charged readings without having to invent a new foundation.
